\section{Обсуждение и выводы}

Полученные результаты разделяют три источника стоимости описания.
Статистическая геометрия определяет байесовскую избыточность и массу
окрестностей малой ошибки. Алгоритм обучения определяет траекторию рисков.
Протокол передачи добавляет цену синхронизации, выбора кандидатов и
восстановления нужного объекта. Один и тот же геометрический показатель
не устраняет различий между этими задачами.

В бернуллиевской модели цена отдельного индекса делает единичные пакеты
неэффективными: каждое обновление требует хотя бы одного бита.
Геометрические расписания ограничивают дивергенцию пакета и число
кандидатов, одновременно уменьшая число сообщений до логарифмического.
Теоремы~\ref{thm:optimal} и~\ref{thm:compute} дают точный ведущий коэффициент
в соответствующем классе кодов. Теорема~\ref{thm:joint} показывает, что
смена внутрипакетного предиктива меняет этот коэффициент даже при
неизменном конечном апостериорном распределении. Поэтому тип предиктива
следует учитывать
вместе со значением оценки сложности.

В факторизованной регрессии один и тот же риск с RLCT $\lambda=1$ допускает
сходящуюся сумму ошибок, степенное накопление и режим $T/\ln^2T$.
Следовательно, при интерпретации накопленной KL как сложности нужно
фиксировать инициализацию, расписание шага и способ усреднения.
Сходство с байесовским коэффициентом на конечном участке траектории само
по себе не доказывает асимптотическое соответствие.

Наконец, для сжимающего ученика история содержит информацию, уже почти не
влияющую на конечный предиктор. Теорема~\ref{thm:suffix} превращает это
наблюдение в конструктивную оценку длины сообщения. Она требует
равномерного сжатия и положительных вероятностей нужных меток; для
произвольной нейросети эти условия автоматически не выполняются.

При сравнении оценок сложности необходимо фиксировать кодируемый объект,
правило декодирования, априорное распределение, порядок наблюдений и режим
обновления. Накопленная KL, сумма верхних оценок, ожидаемая длина индексов
и реализованное сообщение описывают разные стороны одного протокола.

Открытыми остаются точная оптимальная константа для шумного учителя,
сравнение публичных расписаний с адаптивными политиками и анализ
факторизованной регрессии со стохастическим градиентом. Для нейросетей
дополнительные вопросы связаны с устойчивостью декодирования на конечной
точности, стоимостью поиска кандидатов и тем, насколько короткое описание
конечного предиктора сохраняет его поведение на новых данных.
